\documentclass[10pt]{article}
\usepackage[a4paper,margin=0.85in]{geometry}
\usepackage{amsmath,amssymb}
\usepackage{hyperref}
\usepackage{parskip}
\pagestyle{empty}

\newcommand{\eps}{\varepsilon}

\begin{document}

\begin{center}
{\Large\bfseries Contribution Information Sheet}\\[4pt]
{\large VOUCH: Anytime-Valid Certificates for Machine Unlearning\\
via Paired-Canary Hypothesis-Betting}\\[4pt]
Quang-Vinh Dang --- British University Vietnam
\end{center}

\vspace{4pt}
\noindent\textbf{What is the main claim of the paper? Why is this an important
contribution to the machine learning literature?}

The main claim is that machine unlearning can be given a statistical certificate that
is simultaneously (i) exact and finite-sample rather than asymptotic, (ii) valid at
every stopping time an auditor chooses to look, not only at one pre-committed sample
size, and (iii) free of any retrained reference model or distributional assumption on
the model being audited. We obtain this by a single design decision---planting
paired ghost canaries at fine-tuning time, twin records where a fair coin admits one
into training and withholds the other entirely---coupled to a single inferential
engine, betting e-processes from game-theoretic statistics. The coin makes the null
hypothesis "unlearning worked" into an exact, distribution-free statement (twin
exchangeability) that the provider generates rather than assumes, and the betting
engine turns that null into a certificate whose error guarantee survives continuous
monitoring by Ville's inequality.

This matters because verification, not optimisation, is the bottleneck the unlearning
literature has been converging on: dozens of papers now show that behavioural
evaluations of forgetting can be gamed, that quantisation or light fine-tuning
resurfaces "erased" knowledge, and that no unique retrained reference even exists for
the multi-stage pipelines used in practice. Regulators such as the GDPR and the EU AI
Act presume that erasure can be demonstrated, yet the field has had no procedure whose
error rate survives the way audits actually happen---continuously, as deletion
requests stream in, with an auditor free to stop whenever the evidence looks
conclusive. VOUCH is, to our knowledge, the first framework to cast unlearning
certification as sequential equivalence testing and to deliver such a certificate at
LLM scale. It further shows, and we regard this as an important secondary
contribution, that two design choices which look reasonable in isolation---betting a
mixture over multiple membership scores, and composing per-wave certificates across a
deletion stream by multiplying their e-values---are provably \emph{unsound}, and we
give the sound alternatives (per-score consensus; opposite-direction composition for
certificates versus alarms) together with the empirical size of the failure they
prevent (up to 99.6\% false certification for the naive mixture).

\vspace{2pt}
\noindent\textbf{What is the evidence you provide to support your claim? Be precise.}

The paper proves four theorems and two propositions and tests every one of them
empirically, rather than resting the claim on either alone.

\emph{Theoretical evidence.} We prove (i) that under exact unlearning the
paired-canary score difference is an exact fair coin regardless of model, score
function, or data distribution (the exact null); (ii) that certifying with the
per-score minimum e-process controls false certification at $\alpha$ under a
composite null where some scores may still leak, whereas the adaptive-mixture
certificate that seems natural does not; (iii) a matching information-theoretic
lower bound showing the betting engine is within a constant factor of optimal
sample complexity; and (iv) that composing across a stream of deletion waves is
only sound in \emph{opposite} directions for the two claims---certificates by
all-wave consensus, alarms by e-value product---with a counterexample showing why
multiplying certificate e-values across waves is unsound.

\emph{Simulation evidence.} Over 2{,}000 seeds under an adversary that peeks after
every observation and stops at the first crossing, VOUCH's realised type-I error
stays at or below its nominal $\alpha$, while a fixed-sample binomial test used the
same way inflates sevenfold. The same design confirms the composition failure
directly: the mixture-sign certificate falsely certifies 99.6\% of runs against
0.7\% for the per-score minimum, and multiplying per-wave certificate e-values
falsely certifies 96\% of ten-wave histories against 0\% for all-wave consensus.
Median certification time tracks the information-theoretic bound within 15\%.

\emph{End-to-end evidence on real models.} We fine-tune, unlearn (gradient ascent,
GradDiff, negative preference optimisation), and certify real language
models---GPT-2, Pythia-160M, and Phi-1.5, extended to a six-model zoo spanning 0.9M
to 5.1B parameters and training years 2019--2026---on the field's own TOFU and
MUSE-News benchmarks, three seeds each. Detection of un-unlearned models succeeds
in 16 of 18 benchmark runs under the plain sign arm and 18 of 18 under the
magnitude-aware arm we introduce, with zero false certificates issued anywhere;
genuine unlearning is correctly certified in 17 of 18 runs, the one miss resolved
by our own sample-size guidance. Pairing every certificate with a held-out utility
reading exposes a failure invisible to forget-only metrics: gradient ascent earns a
certificate while held-out loss balloons to 58--286 nats against retraining's
1.5--3.4---it "forgets" by destroying the model. Relearning and jailbreak
recoverability probes further revoke two certificates a one-shot evaluation would
have passed. A worked example reproduces one real canary pair verbatim, showing
what each model actually answers and the score gap the certificate consumes.

\vspace{2pt}
\noindent\textbf{What papers by other authors make the most closely related
contributions, and how is your paper related to them?}

Three literatures border this paper, and the closest neighbours in each are treated
explicitly in Section~2 of the manuscript. In statistical verification of unlearning,
Zhang et al.\ (2025) establish that membership-inference scores cannot prove
training membership unless the inclusion of the data was randomised---the soundness
principle our fair-coin canary design is built to satisfy---while Thudi et al.\
(2022) argue that auditable definitions are a prerequisite for any unlearning
guarantee and Zhang et al.\ (2024) show a dishonest provider can defeat existing
verification strategies; Yu et al.\ (2025) prove that path-dependence makes
retrain-equivalence unattainable for multi-stage pipelines, which is why VOUCH's
null is internally generated rather than referenced to a retrained model. The
nearest direct competitor is Pandey et al.\ (2025), which gives a hypothesis-testing
notion of certified unlearning, but theirs is a fixed-sample test tied to a Gaussian
high-dimensional regression regime rather than an anytime-valid, distribution-free
procedure at LLM behavioural scale; certified removal (Guo et al., 2020; Sekhari et
al., 2021) offers rigorous $(\eps_u,\delta_u)$ bounds but only in convex or
near-convex regimes that become vacuous at the scale we target, and we use it as a
semantic anchor (any certified algorithm provably passes VOUCH) rather than as a
competing method. Cryptographic proof-of-unlearning (Eisenhofer et al., 2025)
attests that the correct computation ran, orthogonal to whether forgetting holds
statistically. The inferential engine is borrowed from safe, anytime-valid
inference---e-values, test martingales, and betting confidence sequences (Ramdas et
al., 2023; Gr\"unwald et al., 2024; Waudby-Smith and Ramdas, 2024; Howard et al.,
2021)---whose one prior application adjacent to ours is differential-privacy
auditing (Steinke et al., 2023; Gonz\'alez et al., 2025); that literature
lower-bounds the privacy loss of a training algorithm by mounting a successful
attack, whereas VOUCH upper-bounds residual influence \emph{after} unlearning by
showing no attack in a declared class succeeds---an equivalence test rather than a
significance test. Finally, the unlearning-algorithm and benchmark literature we
certify against---negative preference optimisation (Zhang et al., 2024), the TOFU
(Maini et al., 2024) and MUSE-News (Shi et al., 2025) benchmarks, and the
adversarial-recovery findings of {\L}ucki et al.\ (2025), Zhang et al.\ (2025), and
Deeb and Roger (2024)---is treated throughout as the set of subjects VOUCH
certifies rather than as a competing verification method.

\vspace{2pt}
\noindent\textbf{Have you published parts of your paper before, for instance in a
conference? If so, give details of your previous paper(s) and a precise statement
detailing how your paper provides a significant contribution beyond the previous
paper(s).}

No part of this paper has been published previously in any venue, conference
proceedings, workshop, or preprint server. This submission is the first disclosure
of the VOUCH framework, its theoretical results, and the accompanying experimental
study.

\end{document}
